\chapter{Construction de variables, sélection et importance}\label{ch:ml:feature-engineering-selection-and-importance}

Un modèle utilise deux versions du même signal, calculées sur des fenêtres légèrement différentes, dont la corrélation vaut 0.95. Son
\omterm{def:ml:feature-engineering-selection-and-importance:mdi}{importance par permutation} les classe quatrième et cinquième sur quinze. Réentraîné sans la première, le modèle perd 0.09 point de R-deux,
sans la seconde il en gagne 0.02. Réentraîné sans les deux, il perd 1.05 des 2.67 points qu'il avait. Chaque copie paraît
inutile parce que l'autre la remplace, et un chercheur qui élaguerait les variables «~inutiles~» une à une aurait jeté
deux cinquièmes du modèle. Ce chapitre construit des variables qui se comportent correctement, sélectionne parmi elles et mesure
l'importance avec des méthodes dont on connaît les défaillances à l'avance.

\section{Transformations qui respectent les données}

\begin{definition}[Construction de variables]\label{def:ml:feature-engineering-selection-and-importance:engineering}
La \emph{construction de variables}\index{construction de variables} (feature engineering) désigne l'élaboration des entrées d'un modèle à partir des données brutes~: choisir ce que l'on mesure
(un rendement, un volume, un déséquilibre du carnet d'ordres), sur quelle fenêtre, et comment le transformer (mise à l'échelle, rangs, troncature, différence)
pour que sa relation avec la cible soit stable et que ses valeurs soient connues à l'instant de décision.
\end{definition}

Les variables de marché sont à queue épaisse et non stationnaires. Le livre~7 (chapitre~7) a construit les variables de rendement, de volatilité
et de volume~; le livre~4 a présenté la winsorisation (chapitre~15) et la différenciation fractionnaire (chapitre~17). La transformation compte
davantage que le modèle pour une variable à queue épaisse. Une variable tirée d'une loi de Student $t$ à 1.5 degré de liberté,
agissant sur la cible à travers $0.2\tanh(z)$ (son effet sature), explique 2.17~\% de la variance au mieux. Une régression
linéaire sur la variable brute atteint 0.13~\% hors échantillon (out-of-sample)~: quelques valeurs extrêmes fixent la pente. Winsorisée aux 1\up{er}
et 99\up{e}~centiles, elle atteint 1.27~\%, et ramenée en rangs puis transformée en scores normaux, 1.96~\%. Le classement par date (la
transformation en rangs transversale du livre~7, chapitre~6) est pour la même raison le choix par défaut de ce livre.

\section{Sélection de variables et sélection par stabilité}

\begin{definition}[Sélection de variables, sélection par stabilité]\label{def:ml:feature-engineering-selection-and-importance:selection}
La \emph{sélection de variables}\index{sélection de variables} choisit un sous-ensemble de variables candidates à utiliser. \emph{Sélection par
stabilité}\index{sélection par stabilité} (stability selection) exécute une procédure de sélection (les coefficients non nuls du lasso, les variables les plus importantes d'un modèle)
sur de nombreux demi-échantillons aléatoires, et conserve les variables sélectionnées dans une proportion de ceux-ci supérieure à un seuil
$\pi_{\mathrm{thr}}\in(\frac12, 1]$ (Meinshausen et Bühlmann, 2010).
\end{definition}

\begin{theorem}[Une borne sur les sélections fausses]\label{thm:ml:importance:mb}
Si chaque demi-échantillon sélectionne en moyenne $q$ variables parmi $p$, si les \omterm{def:ml:feature-engineering-selection-and-importance:selection}{sélections des variables} de bruit sont échangeables, et si
la procédure ne fait pas moins bien que deviner au hasard, le nombre espéré de variables de bruit retenues par la \omterm{def:ml:feature-engineering-selection-and-importance:selection}{sélection par stabilité} au seuil
$\pi_{\mathrm{thr}}$ est au plus $q^2/((2\pi_{\mathrm{thr}} - 1)p)$.
\end{theorem}

\admitted{} La démonstration, par une inégalité de type Markov sur les fréquences de sélection, figure dans Meinshausen et Bühlmann (2010).

La tâche du chapitre place les rôles~: parmi quinze variables gaussiennes centrées réduites, $x_1$ agit linéairement, $x_2$ par son carré,
$x_3$ et $x_4$ seulement par leur produit, $x_{1b}$ est presque une copie de $x_1$ (corrélation 0.95, sans effet propre), et dix sont du bruit~;
le signal représente 4~\% de la variance de la cible. Pour la \omterm{def:ml:feature-engineering-selection-and-importance:selection}{sélection par stabilité}, on ajoute 85 colonnes de bruit supplémentaires (100 variables candidates)
et l'on tire 30 demi-échantillons de 5\,000 lignes. Le lasso ne retient que $x_1$ au seuil de 0.6 (il sélectionne 6.7 variables par demi-échantillon~; la borne autorise 2.3 sélections fausses, et aucune ne se produit)~;
le boosting, avec les six variables de plus grand gain, retient $x_1$ et $x_2$ (borne 1.8, aucune fausse). Aucun des deux ne trouve $x_3$
ni $x_4$~: une interaction sans effets principaux est invisible pour un sélecteur linéaire,
et pour un arbre glouton dont la première coupure sur l'une ou l'autre variable ne gagne rien.
La \omterm{def:ml:feature-engineering-selection-and-importance:selection}{sélection par stabilité} contrôle ce qui est sélectionné, non ce qui est manqué.

\section{Importance~: impureté, permutation, Shapley}

\begin{definition}[Diminution moyenne d'impureté, importance par permutation]\label{def:ml:feature-engineering-selection-and-importance:mdi}
La \emph{diminution moyenne d'impureté}\index{diminution moyenne d'impureté} (mean decrease in impurity, MDI, ou importance par gain) d'une variable est la somme des réductions de perte des coupures qui la portent,
sur l'ensemble des arbres, calculée sur les données d'entraînement. L'\emph{importance
par permutation}\index{importance par permutation} (permutation importance) est la baisse d'un score hors échantillon lorsque la colonne de la variable est permutée
au hasard, ce qui rompt son lien avec la cible et avec les autres variables.
\end{definition}

\begin{definition}[Valeur de Shapley, valeur SHAP]\label{def:ml:feature-engineering-selection-and-importance:shapley}
Pour un jeu $v$ sur l'ensemble $P$ de joueurs, la \emph{valeur de Shapley}\index{valeur de Shapley} (Shapley value) du joueur $j$ est
\[
\phi_j = \sum_{S\subseteq P\setminus\{j\}}\frac{|S|!\,(|P| - |S| - 1)!}{|P|!}\bigl(v(S\cup\{j\}) - v(S)\bigr),
\]
l'unique répartition efficace ($\sum_j\phi_j = v(P) - v(\emptyset)$), symétrique, nulle pour un joueur qui n'apporte rien,
et additive sur les jeux. Une \emph{valeur SHAP}\index{valeur SHAP} (SHAP value) est la valeur de Shapley de la variable $j$ pour une prédiction donnée,
où $v(S)$ est la sortie moyenne attendue du modèle lorsque les variables de $S$ sont fixées à leurs valeurs (Lundberg et Lee,
2017)~; pour les ensembles d'arbres, elle se calcule exactement et rapidement (TreeSHAP).
\end{definition}

Les quatre mesures répondent à des questions différentes. La MDI indique où les arbres ont placé leurs coupures sur les données d'entraînement~;
l'\omterm{def:ml:feature-engineering-selection-and-importance:mdi}{importance par permutation} indique de combien le modèle ajusté a besoin d'une colonne sur des données nouvelles~; l'importance par suppression de colonne
(réentraînement sans la variable) indique de combien les \emph{données} en ont besoin~; la moyenne des \omterm{def:ml:feature-engineering-selection-and-importance:shapley}{valeurs SHAP} absolues indique de combien chaque variable déplace les prédictions individuelles du modèle.
Sur la tâche planifiée (la \cref{fig:ml:importance:bars}), des arbres boostés entraînés sur 20\,000 lignes atteignent 2.67~\% sur 20\,000 autres (la vérité atteint 4.01~\%).
La permutation et la suppression de colonne placent $x_2$, $x_3$ et $x_4$ en tête et ne donnent presque rien au bruit. La MDI classe les mêmes quatre variables en tête,
mais attribue à chaque colonne de bruit une part des coupures (la plus grande, 5.4~\% du total, contre 9.5~\% pour $x_1$)~: le gain est acquis sur les données d'entraînement,
bruit compris. SHAP crédite $x_{1b}$ davantage que $x_3$ ou $x_4$~:
il explique ce que fait le modèle, et le modèle utilise la copie.

\begin{omfigure}
\begin{tikzpicture}
\begin{axis}[width=12cm,height=5.8cm,ybar,bar width=5.5pt,ylabel={part du total (\%)},
  symbolic x coords={x1,x2,x3,x4,x1b,noise},xtick=data,xticklabels={$x_1$,$x_2$,$x_3$,$x_4$,$x_{1b}$,bruit max},
  tick label style={font=\small},label style={font=\small},ymin=0,ymax=40,grid=major,grid style={gray!25},
  enlarge x limits=0.1,legend columns=4,legend style={font=\footnotesize,at={(0.5,-0.2)},anchor=north,draw=none,
  fill=none,/tikz/every even column/.append style={column sep=8pt}},table/col sep=comma]
\addplot[fill=gray!50,draw=gray] table[x=feature,y=gain]{figdata/ml/06-feature-engineering-selection-and-importance/importance.csv};
\addplot[fill=omDef!60,draw=omDef] table[x=feature,y=perm]{figdata/ml/06-feature-engineering-selection-and-importance/importance.csv};
\addplot[fill=omThm!55,draw=omThm] table[x=feature,y=drop]{figdata/ml/06-feature-engineering-selection-and-importance/importance.csv};
\addplot[fill=omProp!60,draw=omProp] table[x=feature,y=shap]{figdata/ml/06-feature-engineering-selection-and-importance/importance.csv};
\legend{gain (MDI),permutation,suppression de colonne,moyenne $|$SHAP$|$}
\end{axis}
\end{tikzpicture}

\omcaption{Quatre importances des mêmes arbres boostés sur la tâche planifiée, chacune en part de son total positif~: les
quatre variables de signal, la copie $x_{1b}$, et la plus grande des dix variables de bruit. La permutation, la suppression de colonne et SHAP
sont calculées sur les 20\,000 lignes hors échantillon, le gain sur les lignes d'entraînement. Données~: \texttt{ml\_importance.importances}.}\label{fig:ml:importance:bars}
\end{omfigure}

\section{Substitution et importance par grappes}

\begin{definition}[Effet de substitution, importance des variables par grappes]\label{def:ml:feature-engineering-selection-and-importance:substitution}
L'\emph{effet de substitution}\index{effet de substitution} est la sous-estimation de l'importance d'une variable lorsqu'une variable corrélée peut la remplacer~:
permuter ou supprimer l'une laisse l'autre en place. L'\emph{importance des variables
par grappes}\index{importance des variables par grappes} (clustered feature importance) mesure l'importance de groupes de variables corrélées, permutées ou
supprimées ensemble, les groupes étant obtenus par partitionnement hiérarchique sur $1 - |\text{corrélation}|$ (López de Prado, 2020).
\end{definition}

$x_1$ et sa copie en sont l'exemple. Permutées isolément, elles coûtent 0.76 et 0.41 point de R-deux~; permutées ensemble, 2.72
points, plus que tout le score du modèle, car le modèle a appris à utiliser les deux variables et reçoit désormais du bruit par ces deux canaux.
Supprimées isolément puis réentraînées, elles coûtent 0.09 et $-0.02$ point~; supprimées ensemble, 1.05. Un partitionnement des quinze variables
à une distance de 0.5 trouve exactement un groupe de deux, $\{x_1, x_{1b}\}$, et toutes les autres variables seules~: l'importance par grappes
restitue le signal de $x_1$ à la paire et la désigne ainsi, ce qui est la réponse honnête lorsque les données ne permettent pas de dire
quelle copie compte. Dans les termes de la MDI, la paire partage aussi son crédit (9.5~\% et 8.1~\%), de sorte qu'une règle d'élagage fondée sur le gain aurait hésité
sur les deux.

\section{Quand l'importance trompe}

\begin{method}[Lire les importances]\label{met:ml:importance:reading}
\begin{enumerate}
\item Calculer les importances sur des données hors échantillon (permutation, suppression de colonne)~; lire la MDI calculée sur les données d'entraînement
comme une description de l'ajustement, jamais comme une preuve.
\item Regrouper les variables par corrélation et présenter les importances des grappes~; élaguer des grappes entières.
\item Comparer avec une permutation de la cible (le canari du chapitre~3)~: une variable n'est importante que si elle dépasse
l'importance que reçoivent les variables de bruit.
\item Garder à l'esprit ce que répond chaque mesure~: l'usage que fait le modèle (SHAP, permutation) ou le besoin des données
(suppression de colonne)~; une variable peut être utilisée sans être nécessaire.
\item Chercher séparément les interactions~: un effet passant par un produit peut échapper à tous les criblages marginaux.
\end{enumerate}
\end{method}

L'importance trompe aussi avec le temps. L'importance d'une variable dans un modèle ajusté sur dix ans est une moyenne sur les régimes~; une
variable qui a compté pendant les cinq premières années mais pas pendant les cinq dernières ne garde que la moitié de son ancienne importance et paraît
fiable. Le chapitre~12 mesure l'importance sur des fenêtres glissantes, et le chapitre~21 se demande si une explication est stable
d'un réentraînement à l'autre.

\section{Tutoriel~: quelle variable a compté~?}\label{tut:ml:feature-engineering-selection-and-importance:tut}

\begin{tutorial}
\textbf{Objectif.} Calculer les quatre importances et leurs versions par grappes sur la tâche synthétique, lancer une \omterm{def:ml:feature-engineering-selection-and-importance:selection}{sélection par stabilité}
sur 100 candidats et transformer une variable à queue épaisse de trois manières. \textbf{État final~:} la \cref{fig:ml:importance:bars},
les chiffres du chapitre.
\begin{tutsteps}
\item \textbf{\omterm{def:ml:feature-engineering-selection-and-importance:mdi}{Importance par permutation} et par suppression de colonne}, isolément ou par groupe.
\omcode{code/firm/featimp/firm_featimp.py}{48}{70}{Importances hors échantillon, pour des variables isolées ou des groupes.}\label{lst:ml:importance:perm}
\item \textbf{Grappes et \omterm{def:ml:feature-engineering-selection-and-importance:selection}{sélection par stabilité}.}
\omcode{code/firm/featimp/firm_featimp.py}{77}{103}{Grappes de corrélation, sélection par stabilité et sa borne.}\label{lst:ml:importance:stab}
\item \textbf{Exécuter} \texttt{ml\_importance.importances()}, \texttt{twins()}, \texttt{stability()},
\texttt{heavy\_tail()} et \texttt{fig\_importance.py}.
\end{tutsteps}
\textbf{Pour aller plus loin.} Donner à $x_3$ un petit effet principal et vérifier si la \omterm{def:ml:feature-engineering-selection-and-importance:selection}{sélection par stabilité} trouve désormais
l'interaction~; porter la corrélation de la copie à 0.999 et observer sa \omterm{def:ml:feature-engineering-selection-and-importance:shapley}{valeur SHAP} dépasser celle de $x_1$.
\end{tutorial}

\section{Construction~: importance des variables et sélection}\label{bld:ml:feature-engineering-selection-and-importance:featimp}

\begin{build}
\bfield{Rôle} Chaque rapport de modèle porte des importances hors échantillon, par grappes, et un historique de la manière dont ses variables ont été
sélectionnées.
\bfield{Interface} \texttt{winsorise}, \texttt{rank\_gauss}, \texttt{mdi(model)},
\texttt{permutation\_importance(predict, X, y, score, reps, seed, groups)},
\texttt{drop\_column\_importance(fit\_predict, X, y, Xt, yt, score, groups)}, \texttt{shap\_values(model, X)}, \texttt{cluster\_features(X, threshold)},
\texttt{stability\_selection(select, X, y, n\_sub, frac, seed)}, \texttt{false\_selection\_bound(q, p, threshold)}.
\bfield{Règles} Importances calculées sur des lignes hors échantillon~; les groupes sont permutés conjointement par une même permutation des lignes~; la sélection
est une procédure appelable, si bien que tout modèle peut faire l'objet d'une \omterm{def:ml:feature-engineering-selection-and-importance:selection}{sélection par stabilité}.
\bfield{Tests d'acceptation} \texttt{code/firm/featimp/tests/}~: une variable de pur bruit a une \omterm{def:ml:feature-engineering-selection-and-importance:mdi}{importance par permutation} quasi nulle
et un signal dupliqué n'est retrouvé qu'en tant que groupe~; les \omterm{def:ml:feature-engineering-selection-and-importance:shapley}{valeurs SHAP} somment à la prédiction moins le biais~; le partitionnement place une
quasi-copie avec son original~; la \omterm{def:ml:feature-engineering-selection-and-importance:selection}{sélection par stabilité} conserve une variable forte et écarte le bruit~; la formule de la borne.
\bfield{Pour aller plus loin} Permutation conditionnelle (permuter à l'intérieur des intervalles des variables corrélées)~; importance sur des
fenêtres glissantes~; valeurs d'interaction SHAP pour les paires.
\end{build}

\begin{omsources}
\item N.~Meinshausen et P.~Bühlmann, ``Stability selection'', \emph{Journal of the Royal Statistical Society B} 72(4),
2010.
\item S.~M.~Lundberg et S.-I.~Lee, ``A unified approach to interpreting model predictions'', \emph{NeurIPS}, 2017;
S.~M.~Lundberg et al., ``From local explanations to global understanding with explainable AI for trees'', \emph{Nature
Machine Intelligence} 2, 2020.
\item C.~Strobl, A.-L.~Boulesteix, A.~Zeileis et T.~Hothorn, ``Bias in random forest variable importance measures'',
\emph{BMC Bioinformatics} 8, 2007.
\item M.~López de Prado, \emph{Machine Learning for Asset Managers}, Cambridge University Press, 2020.
\item L.~S.~Shapley, ``A value for $n$-person games'', dans \emph{Contributions to the Theory of Games II}, 1953.
\end{omsources}

\section{Exercices}

\begin{exercise}[$\star$]\label{exo:ml:feature-engineering-selection-and-importance:1}
Une \omterm{def:ml:feature-engineering-selection-and-importance:selection}{sélection par stabilité} sur 200 candidats retient en moyenne 10 variables par demi-échantillon. Calculer la borne aux seuils
0.6 et 0.9.
\end{exercise}

\begin{exercise}[$\star$]\label{exo:ml:feature-engineering-selection-and-importance:2}
Trois joueurs~: $v(\emptyset) = 0$, singletons 1, 1, 0~; paires $\{1,2\} = 3$, $\{1,3\} = 1$, $\{2,3\} = 1$~; la coalition des trois~: 3.
Calculer les \omterm{def:ml:feature-engineering-selection-and-importance:shapley}{valeurs de Shapley}.
\end{exercise}

\begin{exercise}[$\star$]\label{exo:ml:feature-engineering-selection-and-importance:3}
Le R-deux hors échantillon passe de 2.67~\% à 2.58~\% lorsqu'une variable est supprimée et que le modèle est réentraîné. Calculer son importance par suppression de
colonne, et préciser ce qu'on ne peut pas en conclure.
\end{exercise}

\begin{exercise}[$\star\star$]\label{exo:ml:feature-engineering-selection-and-importance:4}
Expliquer pourquoi permuter $x_1$ et $x_{1b}$ ensemble coûte davantage (2.72 points) que tout le R-deux du modèle (2.67~\%).
\end{exercise}

\begin{exercise}[$\star\star$]\label{exo:ml:feature-engineering-selection-and-importance:5}
Expliquer pourquoi la MDI attribue aux variables de bruit jusqu'à 5.4~\% du total, alors que l'\omterm{def:ml:feature-engineering-selection-and-importance:mdi}{importance par permutation} ne leur en accorde que 0.2~\%.
\end{exercise}

\begin{exercise}[$\star\star$]\label{exo:ml:feature-engineering-selection-and-importance:6}
\emph{Trouver la faille.} «~On a éliminé chaque variable dont l'\omterm{def:ml:feature-engineering-selection-and-importance:mdi}{importance par permutation} était inférieure à la médiane, une à la fois,
en réentraînant après chaque suppression, et l'on a conservé ce qui restait.~»
\end{exercise}

\begin{exercise}[$\star\star\star$]\label{exo:ml:feature-engineering-selection-and-importance:7}
\emph{Programmation.} Dans \texttt{ml\_importance.heavy\_tail}, remplacer le $t$ à 1.5 degrés de liberté par celui à 5 degrés de liberté. Relever les
trois R-deux et expliquer pourquoi l'écart entre la variable brute et sa transformation en rangs se referme.
\end{exercise}

\begin{exercise}[$\star\star\star$]\label{exo:ml:feature-engineering-selection-and-importance:8}
Montrer que si $x_{1b} = x_1$ exactement et que la prédiction d'un modèle ne dépend que de $x_1 + x_{1b}$, les \omterm{def:ml:feature-engineering-selection-and-importance:shapley}{valeurs SHAP} des
deux sont égales, et que l'\omterm{def:ml:feature-engineering-selection-and-importance:mdi}{importance par permutation} de la paire n'est pas la somme de leurs importances isolées.
\end{exercise}

\section{Problème~: quelle variable a compté~?}

\begin{problem}[{problème du week-end --- rôles implantés, quatre importances}]\label{pb:ml:feature-engineering-selection-and-importance:1}
La tâche du chapitre~: $x_1$ linéaire, $x_2$ au carré, $x_3x_4$, une copie $x_{1b}$, dix variables de bruit~; des arbres boostés sur 20\,000
lignes, évalués sur 20\,000 autres.

\medskip\noindent\textbf{Partie I --- Transformations.}
\begin{enumerate}
\item Quels scores obtiennent les versions brute, winsorisée et en rangs de la variable à queue épaisse, comparées à la vérité de 2.17~\%~?
\item Pourquoi la transformation en rangs l'emporte-t-elle sur la winsorisation ici~?
\item Quelle transformation le livre utilise-t-il par défaut, et pourquoi~?
\item Quel score le modèle obtient-il, et quelle est la vérité~?
\end{enumerate}

\medskip\noindent\textbf{Partie II --- Quatre importances.}
\begin{enumerate}[resume]
\item Comment chaque mesure classe-t-elle $x_1$ à $x_4$~?
\item Pourquoi SHAP crédite-t-il davantage la copie que $x_3$ ou $x_4$~?
\item Quelle valeur MDI attribue-t-elle à la variable de bruit la plus forte, et pourquoi~?
\item Quelle mesure répond à la question suivante~: «~les données ont-elles besoin de cette variable~?~»
\end{enumerate}

\medskip\noindent\textbf{Partie III --- La copie.}
\begin{enumerate}[resume]
\item Quelle perte de score entraîne la permutation de $x_1$ et $x_{1b}$, isolément puis ensemble~?
\item Quelle perte de score entraîne leur suppression, isolément puis ensemble~?
\item Quel groupe le partitionnement (clustering) identifie-t-il, et à quelle distance~?
\item Qu'aurait donné un élagage une variable à la fois~?
\end{enumerate}

\medskip\noindent\textbf{Partie IV --- Sélection et verdict.}
\begin{enumerate}[resume]
\item Que retiennent le lasso et le boosting avec un seuil de 0.6 sur 100 candidats, et que permet la borne~?
\item Pourquoi les deux méthodes manquent-elles $x_3$ et $x_4$~?
\item Énoncer le \emph{résultat nommé}~: les rangs des variables de signal selon chaque méthode, et le taux de fausses
sélections de la \omterm{def:ml:feature-engineering-selection-and-importance:selection}{sélection par stabilité} au seuil de 0.6.
\item Quelle importance faut-il présenter à un gérant de portefeuille, et laquelle à un comité des risques~?
\item Comment faire en sorte que la \omterm{def:ml:feature-engineering-selection-and-importance:selection}{sélection par stabilité} détecte l'interaction~?
\item Que changerait-il avec des variables qui dérivent dans le temps~?
\item Rédiger la section sur l'importance d'un rapport de modèle en quatre lignes.
\item En une phrase, expliquer ce que mesure une mesure d'importance.
\end{enumerate}
\end{problem}

\section{Questions d'entretien}

\begin{interviewq}[$\star$ \iqroles{chercheur, ingénieur ML}]\label{iq:ml:feature-engineering-selection-and-importance:1}
Quelle est la différence entre l'importance fondée sur l'impureté et l'\omterm{def:ml:feature-engineering-selection-and-importance:mdi}{importance par permutation}, et laquelle jugez-vous fiable~?
\end{interviewq}

\begin{interviewq}[$\star\star$ \iqroles{chercheur}]\label{iq:ml:feature-engineering-selection-and-importance:2}
Deux de vos variables sont corrélées à 95~\% et présentent toutes deux une faible importance. Que faites-vous~?
\end{interviewq}

\begin{interviewq}[$\star\star$ \iqroles{ingénieur ML}]\label{iq:ml:feature-engineering-selection-and-importance:3}
Expliquez les \omterm{def:ml:feature-engineering-selection-and-importance:shapley}{valeurs SHAP} et citez une situation dans laquelle elles induisent en erreur.
\end{interviewq}

\begin{interviewq}[$\star\star$ \iqroles{chercheur}]\label{iq:ml:feature-engineering-selection-and-importance:4}
Comment la \omterm{def:ml:feature-engineering-selection-and-importance:selection}{sélection par stabilité} maîtrise-t-elle les fausses découvertes, et que ne maîtrise-t-elle pas~?
\end{interviewq}

\begin{interviewq}[$\star\star$ \iqroles{chercheur, trader}]\label{iq:ml:feature-engineering-selection-and-importance:5}
Une variable de volume présente parfois des valeurs cent fois supérieures à sa médiane. Comment la fournissez-vous à un modèle~?
\end{interviewq}

\begin{interviewq}[$\star\star\star$ \iqroles{chercheur}]\label{iq:ml:feature-engineering-selection-and-importance:6}
Construisez une cible et deux variables telles que chaque variable, prise seule, ait une corrélation nulle avec la cible, mais
que les deux ensemble la prédisent parfaitement. Que cela implique-t-il pour le criblage des variables~?
\end{interviewq}
